\begin{theorem}[Cauchy–Schwarz Inequality]
For any real numbers $a_{1},\dots ,a_{n}$ and $b_{1},\dots ,b_{n}$,
\[
\Bigl(\sum_{i=1}^{n} a_{i}b_{i}\Bigr)^{2}
\le
\Bigl(\sum_{i=1}^{n} a_{i}^{2}\Bigr)\Bigl(\sum_{i=1}^{n} b_{i}^{2}\Bigr).
\]
Equality holds iff the vectors 
$\mathbf a=(a_{1},\dots ,a_{n})$ and $\mathbf b=(b_{1},\dots ,b_{n})$ are
linearly dependent; i.e. either one of them is the zero vector or there exists
$\lambda\in\mathbb R$ such that $a_{i}=\lambda b_{i}$ for every $i$.
\end{theorem}

\begin{proof}
\textbf{1. Vector notation.}  
Set 
\[
\mathbf a=(a_{1},\dots ,a_{n}),\qquad 
\mathbf b=(b_{1},\dots ,b_{n})\in\mathbb R^{n}.
\]
Their Euclidean inner product and norm are
\[
\mathbf a\cdot\mathbf b=\sum_{i=1}^{n}a_{i}b_{i},\qquad
\|\mathbf a\|^{2}= \sum_{i=1}^{n}a_{i}^{2},\qquad
\|\mathbf b\|^{2}= \sum_{i=1}^{n}b_{i}^{2}.
\]
Thus we must show $(\mathbf a\cdot\mathbf b)^{2}\le\|\mathbf a\|^{2}\,\|\mathbf b\|^{2}$.

\textbf{2. A non‑negative quadratic form.}  
For $t\in\mathbb R$ define
\[
Q(t)=\sum_{i=1}^{n}(a_{i}t-b_{i})^{2}.
\]
Each summand is a square, so $Q(t)\ge 0$ for every $t$.

\textbf{3. Expand $Q(t)$.}
\begin{align*}
Q(t)&=\sum_{i=1}^{n}\bigl(a_{i}^{2}t^{2}-2a_{i}b_{i}t+b_{i}^{2}\bigr)\\
&=\Bigl(\sum_{i=1}^{n}a_{i}^{2}\Bigr)t^{2}
   -2\Bigl(\sum_{i=1}^{n}a_{i}b_{i}\Bigr)t
   +\sum_{i=1}^{n}b_{i}^{2}.
\end{align*}
Hence $Q(t)=At^{2}+Bt+C$ with  
\[
A=\sum_{i=1}^{n}a_{i}^{2},\qquad 
B=-2\sum_{i=1}^{n}a_{i}b_{i},\qquad 
C=\sum_{i=1}^{n}b_{i}^{2}.
\]

\textbf{4. Discriminant condition.}  
A real quadratic polynomial that never becomes negative must have a
non‑positive discriminant, i.e. $B^{2}-4AC\le 0$.

\textbf{5. Compute the discriminant.}
\begin{align*}
B^{2}-4AC
&=\bigl(-2\sum_{i=1}^{n}a_{i}b_{i}\bigr)^{2}
   -4\Bigl(\sum_{i=1}^{n}a_{i}^{2}\Bigr)\Bigl(\sum_{i=1}^{n}b_{i}^{2}\Bigr)\\
&=4\Bigl(\sum_{i=1}^{n}a_{i}b_{i}\Bigr)^{2}
   -4\Bigl(\sum_{i=1}^{n}a_{i}^{2}\Bigr)\Bigl(\sum_{i=1}^{n}b_{i}^{2}\Bigr).
\end{align*}

\textbf{6. Derive the inequality.}  
From $B^{2}-4AC\le 0$ we obtain
\[
4\Bigl(\sum_{i=1}^{n}a_{i}b_{i}\Bigr)^{2}
   -4\Bigl(\sum_{i=1}^{n}a_{i}^{2}\Bigr)\Bigl(\sum_{i=1}^{n}b_{i}^{2}\Bigr)\le 0,
\]
and dividing by the positive constant $4$ yields
\[
\Bigl(\sum_{i=1}^{n} a_{i}b_{i}\Bigr)^{2}
\le
\Bigl(\sum_{i=1}^{n} a_{i}^{2}\Bigr)\Bigl(\sum_{i=1}^{n} b_{i}^{2}\Bigr),
\]
which is the Cauchy–Schwarz inequality.

\textbf{7. Equality condition.}  
Equality holds exactly when the discriminant is zero, i.e. when
\[
\sum_{i=1}^{n}(a_{i}t-b_{i})^{2}=0
\quad\text{for }t_{0}=-\frac{B}{2A}
   =\frac{\sum_{i=1}^{n}a_{i}b_{i}}{\sum_{i=1}^{n}a_{i}^{2}}.
\]
The sum of squares can vanish only if each term vanishes, so
$a_{i}t_{0}=b_{i}$ for all $i$. Hence either $\mathbf a=\mathbf 0$ (trivial
equality) or there exists $\lambda=t_{0}$ such that $a_{i}=\lambda b_{i}$ for every $i$.
Thus equality occurs iff $\mathbf a$ and $\mathbf b$ are linearly dependent
(including the degenerate zero‑vector case).

\end{proof}